\textbf{19.} Докажите $\mathbf{NL}$-полноту языка ориентированных ациклических графов.

\textbf{Решение.} Как известно, $\mathbf{NL}=\mathbf{coNL}$, отчего достаточно перейти к дополнению языка из условия и показать $\mathbf{NL}$-полноту языка ориентированных графов с циклом.

Принадлежность $\mathbf{NL}$ очевидна: достаточно в качестве сертификата дать вершины цикла в порядке обхода, а верификатор будет идти по последовательности слева направо (напоминаем, что в сертификатном определении $\mathbf{NL}$ читать второй аргумент можно только так) и проверять, соединены ли они ребром, а потом он убедится, что первая и последняя вершина совпадают (первую вершину можно куда-нибудь записать на рабочую ленту).

С $\mathbf{NL}$-трудностью всё не так просто. Будем доказывать её с помощью сведения языка $\mathsf{PATH}$ (известно, что он $\mathbf{NL}$-трудный) к данному. Пусть нам дан граф $G$ с вершинами $V$ в количестве $n$ штук, и вершины $s$ и $t$, наличие пути между которыми надо проверить. Построим следующий слоистый граф:
\begin{center}
    \begin{asy}
import graph;
size(220,0);

unitsize(2cm);
        

draw((1,1)--(-1,1)--(-1,-1)--(1,-1)--cycle);
label("$V$", (0, 0));
label("$u$", (-0.3, 0.7), SW);
dot((-0.3, 0.7), filltype=FillDraw(drawpen=black));
dot((0.6, -0.2), filltype=FillDraw(drawpen=black));
label("$v$", (0.6, -0.2), SW);
draw((-0.3, 0.7)--(0.6, -0.2), EndArrow);

label("$\Longrightarrow$", (1.5, 0));


draw((2,1)--(2.5,1)--(2.5,-1)--(2,-1)--cycle);
label("$V_1$", (2.25, -1), N);
draw((3,1)--(3.5,1)--(3.5,-1)--(3,-1)--cycle);
label("$V_2$", (3.25, -1), N);
draw((5,1)--(5.5,1)--(5.5,-1)--(5,-1)--cycle);
label("$V_{n}$", (5.25, -1), N);
label("$\ldots$", (4.25, 0));

dot((2.25, 0.7), filltype=FillDraw(drawpen=black));
label("$u$", (2.25, 0.7), N);
dot((2.25, -0.2), filltype=FillDraw(drawpen=black));
label("$v$", (2.25, -0.2), S);
dot((3.25, 0.7), filltype=FillDraw(drawpen=black));
label("$u$", (3.25, 0.7), N);
dot((3.25, -0.2), filltype=FillDraw(drawpen=black));
label("$v$", (3.25, -0.2), S);
dot((5.25, 0.7), filltype=FillDraw(drawpen=black));
label("$u$", (5.25, 0.7), N);
dot((5.25, -0.2), filltype=FillDraw(drawpen=black));
label("$v$", (5.25, -0.2), S);

draw((2.25, 0.7)--(3.25, -0.2), EndArrow);
draw((2.25, 0.7)--(3.25, 0.7), EndArrow);
draw((2.25, -0.2)--(3.25, -0.2), EndArrow);

draw((3.25, 0.7)--(4.25, -0.2), EndArrow);
draw((3.25, 0.7)--(4.25, 0.7), EndArrow);
draw((3.25, -0.2)--(4.25, -0.2), EndArrow);

draw((4.25, 0.7)--(5.25, -0.2), EndArrow);
draw((4.25, 0.7)--(5.25, 0.7), EndArrow);
draw((4.25, -0.2)--(5.25, -0.2), EndArrow);

dot((5.25, 0.4), filltype=FillDraw(drawpen=red));
label("$t$", (5.25, 0.4), SW);
dot((2.25, 0.4), filltype=FillDraw(drawpen=red));
label("$s$", (2.25, 0.4), SE);
draw((5.25, 0.4){right}..(6, 1.5){left}..{left}(1.5, 1.5)..{right}(2.25, 0.4), red, EndArrow);
label("$e$", (3.75, 1.5), S, red);


    \end{asy}
\end{center}

Мы дублируем множество вершин $n$ раз, получится $n$ слоёв, и внутри каждого рёбер нет. Если в исходном графе было ребро $(u, v)$, то между соседними слоями $V_i$ и $V_{i+1}$ мы проводим ребро из $u$, которая лежит в $V_i$, в вершину $v$, которая лежит в $V_{i+1}$ (см. картинку). Также проведём рёбра между одноимёнными вершинами соседних слоёв слева направо. Последний штрих --- проводим ребро $e$ из вершины $t$ последнего слоя $V_{n}$ и вершину $s$ первого слоя $V_1$. 

Пожинаем плоды нашей умственной деятельности. Если в исходном графе был путь из $s$ в $t$, то его можно пройти и в новом графе, начиная с вершины первого слоя и проходя вправо. Количество слоёв гарантирует, что мы не упрёмся в $V_n$ раньше времени. Когда дойдём до $t$, переходим по рёбрам между $t$-шками, и при достижении последнего слоя переходим в начало по ребру $e$ --- получился цикл. Обратно, если в нашем графе и есть цикл, то он обязан содержать ребро $e$, ведь все остальные рёбра идут слева направо, в порядке увеличения номера слоя, так что цикл из них получить не выйдет. Но тогда в нашем графе есть путь из первой $s$ в последнюю $t$ (возможно, использующий рёбра между одноимёнными вершинами). Значит он же есть и в исходном графе.

Отметим также, что полученное сведение является логарифмическим. Если считать, что граф задан списком рёбер, то мы можем при считывании очередного ребра $(u, v)$ пройтись циклом по $i=1, \ldots, n-1$, выводя рёбра из $u_i$ в $v_{i+1}$. Хранение счётчика требует лишь логарифм памяти.